\begin{table}[t]
\centering
\footnotesize
\caption{\model evaluates \num API-based or OSS LLMs on LLM-as-Agent challenges. Models annotated with * are evaluated after task weights are computed.}
\vspace{-2mm}
\renewcommand\tabcolsep{0.7pt}
\renewcommand\arraystretch{1.1}
\label{tab:model_details}
\resizebox{0.95\columnwidth}{!}{
\begin{tabular}{@{}lcccclcccc@{}}
\toprule
% \rowcolor[HTML]{ECF4FF}
Model                                                                                     & \#Size  &  Form    &  Ver.            &   Creator                                                                                            &
Model                                                                                     & \#Size  &  Form    &  Ver.            &   Creator                                                                                            \\
\midrule

\cellcolor[HTML]{ECF4FF}  \texttt{gpt-4}~\citep{openai2023gpt}                 & \cellcolor[HTML]{ECF4FF}  N/A & \cellcolor[HTML]{ECF4FF}  api  & \cellcolor[HTML]{ECF4FF}  0613     & \cellcolor[HTML]{ECF4FF}                                                                                 & \cellcolor[HTML]{FFFFFF}  \texttt{dolly-12b}~\citep{DatabricksBlog2023DollyV2} & \cellcolor[HTML]{FFFFFF}  12B & \cellcolor[HTML]{FFFFFF}  open & \cellcolor[HTML]{FFFFFF}  v2    & \cellcolor[HTML]{FFFFFF}  Databricks                  \\
\cellcolor[HTML]{ECF4FF}  \texttt{gpt-3.5-turbo}~\citep{ChatGPT}               & \cellcolor[HTML]{ECF4FF}  N/A & \cellcolor[HTML]{ECF4FF}  api  & \cellcolor[HTML]{ECF4FF}  0613     & \cellcolor[HTML]{ECF4FF}                                                                                 & \cellcolor[HTML]{ECF4FF}  \texttt{llama2-70b}~\citep{touvron2023llama}         & \cellcolor[HTML]{ECF4FF}  70B & \cellcolor[HTML]{ECF4FF}  open & \cellcolor[HTML]{ECF4FF}  chat  & \cellcolor[HTML]{ECF4FF}                              \\
\cellcolor[HTML]{ECF4FF}  \texttt{text-davinci-003}~\citep{ouyang2022training} & \cellcolor[HTML]{ECF4FF}  N/A & \cellcolor[HTML]{ECF4FF}  api  & \cellcolor[HTML]{ECF4FF}  -        & \cellcolor[HTML]{ECF4FF}                                                                                 & \cellcolor[HTML]{ECF4FF}  \texttt{llama2-13b}~\citep{touvron2023llama}         & \cellcolor[HTML]{ECF4FF}  13B & \cellcolor[HTML]{ECF4FF}  open & \cellcolor[HTML]{ECF4FF}  chat  & \cellcolor[HTML]{ECF4FF}                              \\
\cellcolor[HTML]{ECF4FF}  \texttt{text-davinci-002}~\citep{ouyang2022training} & \cellcolor[HTML]{ECF4FF}  N/A & \cellcolor[HTML]{ECF4FF}  api  & \cellcolor[HTML]{ECF4FF}  -        & \cellcolor[HTML]{ECF4FF}  \multirow{-4}{*}{OpenAI}                                                       & \cellcolor[HTML]{ECF4FF}  \texttt{llama2-7b}~\citep{touvron2023llama}          & \cellcolor[HTML]{ECF4FF}  7B  & \cellcolor[HTML]{ECF4FF}  open & \cellcolor[HTML]{ECF4FF}  chat  & \cellcolor[HTML]{ECF4FF}  \multirow{-3}{*}{Meta}      \\
\cellcolor[HTML]{FFFFFF}  \texttt{claude-3*}~\citep{Claude-3}                  & \cellcolor[HTML]{FFFFFF}  N/A & \cellcolor[HTML]{FFFFFF}  api  & \cellcolor[HTML]{FFFFFF}  opus     & \cellcolor[HTML]{FFFFFF}                                                                                 & \cellcolor[HTML]{FFFFFF}  \texttt{guanaco-65b}~\citep{dettmers2023qlora}       & \cellcolor[HTML]{FFFFFF}  65B & \cellcolor[HTML]{FFFFFF}  open & \cellcolor[HTML]{FFFFFF}  -     & \cellcolor[HTML]{FFFFFF}                              \\
\cellcolor[HTML]{FFFFFF}  \texttt{claude-2}~\citep{Claude-2}                   & \cellcolor[HTML]{FFFFFF}  N/A & \cellcolor[HTML]{FFFFFF}  api  & \cellcolor[HTML]{FFFFFF}  -        & \cellcolor[HTML]{FFFFFF}                                                                                 & \cellcolor[HTML]{FFFFFF}  \texttt{guanaco-33b}~\citep{dettmers2023qlora}       & \cellcolor[HTML]{FFFFFF}  33B & \cellcolor[HTML]{FFFFFF}  open & \cellcolor[HTML]{FFFFFF}  -     & \cellcolor[HTML]{FFFFFF}  \multirow{-2}{*}{Meta}      \\
\cellcolor[HTML]{FFFFFF}  \texttt{claude}~\citep{Claude}                       & \cellcolor[HTML]{FFFFFF}  N/A & \cellcolor[HTML]{FFFFFF}  api  & \cellcolor[HTML]{FFFFFF}  v1.3     & \cellcolor[HTML]{FFFFFF}                                                                                 & \cellcolor[HTML]{ECF4FF}  \texttt{vicuna-33b}~\citep{chiang2023vicuna}         & \cellcolor[HTML]{ECF4FF}  33B & \cellcolor[HTML]{ECF4FF}  open & \cellcolor[HTML]{ECF4FF}  v1.3  & \cellcolor[HTML]{ECF4FF}                              \\
\cellcolor[HTML]{FFFFFF}  \texttt{claude-instant}~\citep{Claude}               & \cellcolor[HTML]{FFFFFF}  N/A & \cellcolor[HTML]{FFFFFF}  api  & \cellcolor[HTML]{FFFFFF}  v1.1     & \cellcolor[HTML]{FFFFFF}  \multirow{-4}{*}{Anthropic}                                                    & \cellcolor[HTML]{ECF4FF}  \texttt{vicuna-13b}~\citep{chiang2023vicuna}         & \cellcolor[HTML]{ECF4FF}  13B & \cellcolor[HTML]{ECF4FF}  open & \cellcolor[HTML]{ECF4FF}  v1.5  & \cellcolor[HTML]{ECF4FF}                              \\
\cellcolor[HTML]{ECF4FF}  \texttt{chat-bison-001}~\citep{anil2023palm}         & \cellcolor[HTML]{ECF4FF}  N/A & \cellcolor[HTML]{ECF4FF}  api  & \cellcolor[HTML]{ECF4FF}  -        & \cellcolor[HTML]{ECF4FF}  Google                                                                         & \cellcolor[HTML]{ECF4FF}  \texttt{vicuna-7b}~\citep{chiang2023vicuna}          & \cellcolor[HTML]{ECF4FF}  7B  & \cellcolor[HTML]{ECF4FF}  open & \cellcolor[HTML]{ECF4FF}  v1.5  & \cellcolor[HTML]{ECF4FF}  \multirow{-3}{*}{LMSYS}     \\
\cellcolor[HTML]{FFFFFF}  \texttt{glm-4*}~\citep{zeng2022glm,du2022glm}        & \cellcolor[HTML]{FFFFFF}  N/A & \cellcolor[HTML]{FFFFFF}  api  & \cellcolor[HTML]{FFFFFF}  -        & \cellcolor[HTML]{FFFFFF}                                                                                 & \cellcolor[HTML]{FFFFFF}  \texttt{openchat-13b}~\citep{openchat}               & \cellcolor[HTML]{FFFFFF}  13B & \cellcolor[HTML]{FFFFFF}  open & \cellcolor[HTML]{FFFFFF}  v3.2  & \cellcolor[HTML]{FFFFFF}  Tsinghua                    \\
\cellcolor[HTML]{FFFFFF}  \texttt{chatglm-6b}~\citep{zeng2022glm,du2022glm}    & \cellcolor[HTML]{FFFFFF}  6B  & \cellcolor[HTML]{FFFFFF}  open & \cellcolor[HTML]{FFFFFF}  v1.1     & \cellcolor[HTML]{FFFFFF}                                                                                 & \cellcolor[HTML]{ECF4FF}  \texttt{wizardlm-30b}~\citep{xu2023wizardlm}         & \cellcolor[HTML]{ECF4FF}  30B & \cellcolor[HTML]{ECF4FF}  open & \cellcolor[HTML]{ECF4FF}  v1.0  & \cellcolor[HTML]{ECF4FF}                              \\
\cellcolor[HTML]{FFFFFF}  \texttt{codegeex2-6b}~\citep{zheng2023codegeex}      & \cellcolor[HTML]{FFFFFF}  6B  & \cellcolor[HTML]{FFFFFF}  open & \cellcolor[HTML]{FFFFFF}  -        & \cellcolor[HTML]{FFFFFF}  \multirow{-3}{*}{\begin{tabular}[c]{@{}c@{}}Tsinghua \\ \& Zhipu\end{tabular}} & \cellcolor[HTML]{ECF4FF}  \texttt{wizardlm-13b}~\citep{xu2023wizardlm}         & \cellcolor[HTML]{ECF4FF}  13B & \cellcolor[HTML]{ECF4FF}  open & \cellcolor[HTML]{ECF4FF}  v1.0  & \cellcolor[HTML]{ECF4FF}  \multirow{-2}{*}{Microsoft} \\
\cellcolor[HTML]{ECF4FF}  \texttt{codellama-34b}~\citep{roziere2023code}       & \cellcolor[HTML]{ECF4FF}  34B & \cellcolor[HTML]{ECF4FF}  open & \cellcolor[HTML]{ECF4FF}  instruct & \cellcolor[HTML]{ECF4FF}                                                                                 & \cellcolor[HTML]{FFFFFF}  \texttt{koala-13b}~\citep{geng2023koala}             & \cellcolor[HTML]{FFFFFF}  13B & \cellcolor[HTML]{FFFFFF}  open & \cellcolor[HTML]{FFFFFF}  -     & \cellcolor[HTML]{FFFFFF}  UCB                         \\
\cellcolor[HTML]{ECF4FF}  \texttt{codellama-13b}~\citep{roziere2023code}       & \cellcolor[HTML]{ECF4FF}  13B & \cellcolor[HTML]{ECF4FF}  open & \cellcolor[HTML]{ECF4FF}  instruct & \cellcolor[HTML]{ECF4FF}                                                                                 & \cellcolor[HTML]{ECF4FF}  \texttt{oasst-12b}~\citep{2023openassistant}         & \cellcolor[HTML]{ECF4FF}  12B & \cellcolor[HTML]{ECF4FF}  open & \cellcolor[HTML]{ECF4FF}  sft-4 & \cellcolor[HTML]{ECF4FF}  LAION                       \\
\cellcolor[HTML]{ECF4FF}  \texttt{codellama-7b}~\citep{roziere2023code}        & \cellcolor[HTML]{ECF4FF}  7B  & \cellcolor[HTML]{ECF4FF}  open & \cellcolor[HTML]{ECF4FF}  instruct & \cellcolor[HTML]{ECF4FF}  \multirow{-3}{*}{Meta}                                                         &                                                                                &                               &                                &                                 &                                                       \\ 
\bottomrule
\end{tabular}
}
\vspace{-7mm}
\end{table}

\section{LLM-as-Agent: Definition and Preliminary} \label{sec:preliminary}
Here, we formalize the terms for describing the evaluation of LLMs as agents and the necessary preliminary knowledge for using LLMs in the context of agent evaluation.

\vpara{Definition: Interactive Evaluation of LLM-as-Agent.}
The interactive evaluation of LLM-as-Agent could be regarded as a Partially Observable Markov Decision Process ($\mathcal{S}, \mathcal{A}, \mathcal{T}, \mathcal{R}, \mathcal{U}, \mathcal{O}$), which comprises state space $\mathcal{S}$, action space $\mathcal{A}$, transition function $\mathcal{T}: \mathcal{S}\times\mathcal{A}\to\mathcal{S}$, reward assigning function $\mathcal{R}$, task instruction space $\mathcal{U}$, and observation space $\mathcal{O}$.
Here, we denote an LLM agent as $\mathcal{M}$.

\vpara{Chain-of-Thought (CoT) and Other Reasoning Strategies.}
Since LLM-as-Agent requires LLMs' strong reasoning ability, CoT~\citep{wei2022chain}, which has been considered a de facto strategy in related evaluation together with actions~\citep{yao2023react}, is also adopted in \model.
Despite many improved strategies proposed later, such as introducing ensemble~\citep{wang2023self}, reflection~\citep{shinn2023reflexion}, and search~\citep{yao2023tree}, we evaluate LLMs with the most primitive CoT in \model.
Without multiple trials, repeated generations, or complicated strategies, CoT is the easiest, cheapest, and most common way for people to deploy LLM agents.

\vpara{Typical Types of Finish Reasons.}
Despite LLMs' capabilities, we show in \model that even the strongest \texttt{gpt-4} is not qualified as a practically usable agent.
We identify and categorize finish reasons of LLM agents on \model tasks into five typical types:
\begin{itemize}[leftmargin=1.5em,itemsep=0pt,parsep=0.2em,topsep=0.0em,partopsep=0.0em]
\item \textbf{Context Limit Exceeded (CLE)}: the length of interaction history exceeds the LLM's maximum context length (only happened in 2,048-length LLMs \texttt{text-davinci-002} and \texttt{003}).
\item \textbf{Invalid Format (IF)}: the agent does not follow the format instruction.
\item \textbf{Invalid Action (IA)}: the agent follows the format instruction, but its selected action is invalid.
\item \textbf{Task Limit Exceeded (TLE)}: the agent does not solve the problem after reaching the predefined maximum interaction rounds or begins to do repeated generations for many rounds.
\end{itemize}
and Complete (task ends normally).
While IF and IA are mostly caused by LLMs' poor instruction following, TLE often indicates a weak multi-turn ability in certain tasks.